\scanpage{57}
\mspage{24}
de l'ens à $N(x)+1$ éléments, où $N(x)=\operatorname{card}(\Sigma_{<x})$.
Ou encore, si on regarde l'application-but
\[
\Sigma^{<}(d)\longrightarrow\unsure{\Sigma}=\Sigma^{<}(0)
\]
\[
(x_0<\cdots<x_d)\longmapsto x_d
\]
la fibre de $x\in\Sigma$ s'identifie à l'ens \deleted{des parties}
$\mathfrak P_d(\{1,\ldots,N(x)\})$.
La structure relative de $\Sigma^{<}(1)$ sur $\Sigma=\Sigma^{<}(0)$
est donc celle d'un ens tot. ordonné relatif, qui
permet de reconstituer le système des $\Sigma^{<}(d)$
\deleted{avec} $(d\in\mathbb N)$ avec les applications semi-simpliciales
correspondant aux morphismes cofinaux entre
les ens tot. ordonnés-types.
\margin{\unsure{Tout ceci est donc à faire indépendt de toute condition topologique} \unsure{\textcircled{a}}.}

\textcircled{c} $\forall\ x\in\Sigma$, les germes \inserted{$B_\alpha$} de sous-espaces \inserted{bordants} de $\Sigma$ en $x$
correspondants, par \textcircled{a}, aux $z_\alpha\in\Sigma^{<}(1)$ au dessus
de $x$ (par $\Sigma^{<}(1)\xrightarrow{b}\Sigma$) sont mutuellement transver-
saux (et à fortiori distincts).
Les germes intersections de certains des $B_\alpha$
correspondent donc aux pts des $\Sigma^{<}(d)$ au dessus
de $x$ (par $\Sigma^{<}(d)\overset{b}{\unsure{-|\longrightarrow}}\Sigma$), comme
images du \deleted{\unsure{dit}} germe de $\Sigma^{<}(d)$ en ces pts.
Les espaces $\unsure{\Sigma^{*}(d)|\Sigma^{<}(d)}$ $(d\in\mathbb N)$ et les
applications $\unsure{\tfrac12}$-simpliciales cofinales entre eux, \deleted{\unsure{satisf}}
ou encore l'espace relatif $\Sigma^{<}(1)$ sur $\Sigma^{<}(0)$ et la
structure d'ordre relative, est décrit entièrement
\mspage{56}
